% ============================================================
%  AUGMENTATION 2 -- Feature design rationale
%  Insert this paragraph in Section 2.6 between:
%    end of \paragraph{Features.} block (after the pair-feature equation eq:pairfeature)
%  and the start of:
%    \paragraph{Architecture.}
% ============================================================

\paragraph{Feature design rationale.}
The 7-dimensional self-feature and 6-dimensional pair feature are organised as three complementary information channels that the scorer needs in order to discriminate Dirichlet side-lobe atoms from noise-driven candidates. Self-features 1, 2, and 5 expose the candidate's data-driven correlation strength against three different normalisations: the absolute matched-filter magnitude $\log|\mbPhi_i^H \mathbf{r}^{(t)}|$, the margin against the weakest active atom $\log|\mbPhi_i^H \mathbf{r}^{(t)}| - \min_{s \in \hat{S}} \log|\hat{h}_s|$ that places the candidate on the same logarithmic axis as the active support's tail, and the residual-normalised projection $\|\mbPhi_i^H \mathbf{r}^{(t)}\|/\|\mathbf{r}^{(t)}\|$ that controls for the overall residual energy at iteration $t$. Self-features 3 and 4, the normalised delay and Doppler indices $l_i / \lmax$ and $(k_i + \kmax)/(2\kmax)$, locate the candidate on the physical lattice and let the scorer learn position-dependent priors --- for instance, that bins near $\pm \kmax$ are more likely to host fractional-leakage residuals than bins near zero. Self-feature 6, the candidate's rank among the $M = 96$ retained atoms, provides a tie-breaker that lets the scorer condition on whether a candidate sits near the top of the matched-filter shortlist or at its tail. Self-feature 7, the binary indicator that $(l_i, k_i)$ lies within $\kleak$ Doppler bins of an active atom, is the most directly mechanism-coupled feature: it flags the geometric signature of a Dirichlet side-lobe neighbour and is the feature that most cleanly distinguishes the high-SNR exploitable regime from the low-SNR noise-dominated regime developed in \S\ref{ssec:disc_head}.

The 6-dimensional pair feature $\mathbf{v}_{i,s}$ extends the same reasoning to the relational structure between candidate $i$ and each active atom $s$. The integer offsets $(l_i - l_s, k_i - k_s)$ are the direct encoding of the $(\Delta l, \Delta k) = (0, \pm 1)$ Dirichlet side-lobe signature; combined with the active atom's log-gain $\log|\hat{h}_s|$ and its gain prominence $|\hat{h}_s| / \max_{s'}|\hat{h}_{s'}|$, they let the scorer learn that a candidate at $(0, \pm 1)$ offset from a high-gain active atom is far more likely to be a side-lobe than the same candidate at the same offset from a weak active atom. The pair correlation $|\mbPhi_i^H \mbPhi_s|$ measures dictionary-column similarity directly and supplements the integer-offset signal with a continuous coherence measure, so that pairs of atoms whose chirp-domain responses overlap geometrically can be flagged even when the integer-offset pattern is ambiguous. The normalised Euclidean distance $\|(l_i - l_s,\, k_i - k_s)\|_2 / (\lmax + 2\kmax)$ provides a single scalar that the order-invariant pooling of the architecture (described next) can aggregate across all $|\hat{S}|$ active atoms without losing the geometric information. Taken together, the feature design encodes the rule that a candidate should be promoted when it is a low-Hamming-distance Doppler neighbour of an active atom of non-trivial gain and when its matched-filter response sits above the active-set tail; this is the rule that the oracle-headroom analysis of \S\ref{ssec:probes} validates as covering the bulk of the achievable support-refinement gap, and it is the rule that the scorer is trained to recover.
